\section{SGD 在期望中下降但噪声会留下平台}
\label{sec:09-Convex-Optimization/08-Nonconvex-and-Stochastic-Optimization/04}

训练跑到第两万步，损失在 \(0.42\) 附近上下抖 \(0.03\)，再跑两万步还是这样。按上一篇的口径，你会开始怀疑地形——是不是掉进了某个平台。但在动手做曲率探测之前，先做一件更便宜的事：把学习率从 \(3\times10^{-4}\) 降到 \(3\times10^{-5}\)。曲线立刻往下掉了一截，稳定在 \(0.31\)，然后继续横着抖。

这个阶梯形状在训练日志里太常见，常见到容易被当成理所当然。它值得追问：如果 \(0.42\) 那个高度是模型的能力上限，降学习率不该有任何效果；既然有效果，那个高度就不是模型的性质，而是学习率的性质。

\subsection{无偏消不掉方差}

设目标是 \(F(\theta)=\E_\xi[\ell(\theta;\xi)]\)，每步只能拿到一个 mini-batch 上的随机梯度 \(g_t\)，满足条件无偏

\[
\E[g_t\given\theta_t]=\nabla F(\theta_t),\qquad
\E\bigl[\norm{g_t-\nabla F(\theta_t)}^2\given\theta_t\bigr]\le\sigma^2 .
\]

这两条里藏着一个不对称，是本篇后面所有结论的来源。第一条说的是偏差为零，而偏差本来就随接近最优而衰减——真梯度自己会趋于零。第二条里的 \(\sigma^2\) 不衰减。走到最优点附近时，\(\nabla F\) 已经很小，\(g_t\) 却仍然有 \(\sigma\) 那么长；更新的步长 \(\eta\norm{g_t}\) 不会随之缩短。全批量梯度下降是自己刹车的，随机梯度下降不是。

\subsection{带噪的下降不等式多出一项常数}

设 \(F\) 是 \(L\)-光滑的，更新为 \(\theta_{t+1}=\theta_t-\eta g_t\)。下降引理照样适用（它只需要光滑性，与用的是不是真梯度无关）：

\[
F(\theta_{t+1})\le F(\theta_t)-\eta\,\nabla F(\theta_t)^{\trans}g_t+\frac{L\eta^2}{2}\norm{g_t}^2 .
\]

给定 \(\theta_t\) 取条件期望。中间那项由无偏性变成 \(-\eta\norm{\nabla F(\theta_t)}^2\)；末项用方差分解 \(\E[\norm{g_t}^2\given\theta_t]=\norm{\nabla F(\theta_t)}^2+\E[\norm{g_t-\nabla F(\theta_t)}^2\given\theta_t]\) 拆开，得到

\[
\E[F(\theta_{t+1})\given\theta_t]
\le F(\theta_t)-\eta\Bigl(1-\frac{L\eta}{2}\Bigr)\norm{\nabla F(\theta_t)}^2
+\frac{L\eta^2\sigma^2}{2}.
\]

再设 \(F\ge F_{\inf}\)，取 \(0<\eta\le1/L\)，对上式取全期望并从 \(t=0\) 累加到 \(T-1\)：

\[
\frac1T\sum_{t=0}^{T-1}\E\norm{\nabla F(\theta_t)}^2
\le\frac{2\bigl(F(\theta_0)-F_{\inf}\bigr)}{\eta T}+L\eta\sigma^2 .
\]

条件对账。用到光滑性、下界、无偏性、方差有界，以及步长为常数（否则累加时 \(\eta\) 提不出来）。没有用到凸性，没有用到 PL，也没有用到噪声的分布形状——只用了它的二阶矩。

右端两项的行为完全不同。第一项按 \(1/T\) 消失，是确定性部分留下的痕迹。第二项 \(L\eta\sigma^2\) 与 \(T\) 无关，跑多久都在。它就是那个平台：固定步长的 SGD 不会让梯度范数趋于零，只会把它压进一个半径由 \(\eta\) 和 \(\sigma\) 共同标定的邻域，然后在里面转。

平台高度正比于 \(\eta\)。学习率降到十分之一，平台降到十分之一——开头那次从 \(0.42\) 掉到 \(0.31\) 的动作，读的就是这个式子。

\begin{center}
\begin{tikzpicture}[x=0.85cm,y=1.05cm,font=\small,>=stealth]
  \draw[->,black!55] (0,0) -- (10.6,0) node[right,font=\footnotesize,black!70] {步数};
  \draw[->,black!55] (0,0) -- (0,3.25) node[above,font=\footnotesize,black!70] {训练损失};
  \draw[black!30,dashed] (0,1.07) -- (10.2,1.07);
  \draw[black!30,dashed] (3.35,0) -- (3.35,2.55);
  \draw[black!30,dashed] (6.35,0) -- (6.35,1.9);
  \draw[black!40,thick] plot coordinates {
    (0,2.90) (0.4,2.35) (0.8,1.95) (1.2,1.68) (1.6,1.48) (2.0,1.34)
    (2.4,1.24) (2.8,1.17) (3.2,1.10) (3.5,1.14) (3.8,1.06) (4.1,1.12)
    (4.4,1.04) (4.7,1.11) (5.0,1.03) (5.3,1.10) (5.6,1.05) (5.9,1.12)
    (6.2,1.04) (6.5,1.09) (6.8,1.03) (7.1,1.11) (7.4,1.05) (7.7,1.10)
    (8.0,1.02) (8.3,1.09) (8.6,1.04) (8.9,1.11) (9.2,1.03) (9.5,1.08)
    (9.8,1.05)};
  \draw[black!80,thick] plot coordinates {
    (0,2.90) (0.4,2.35) (0.8,1.95) (1.2,1.68) (1.6,1.48) (2.0,1.34)
    (2.4,1.24) (2.8,1.17) (3.2,1.10) (3.5,0.92) (3.8,0.80) (4.1,0.72)
    (4.4,0.68) (4.7,0.71) (5.0,0.66) (5.3,0.70) (5.6,0.65) (5.9,0.69)
    (6.2,0.66) (6.5,0.55) (6.8,0.47) (7.1,0.42) (7.4,0.40) (7.7,0.42)
    (8.0,0.39) (8.3,0.41) (8.6,0.39) (8.9,0.41) (9.2,0.39) (9.5,0.40)
    (9.8,0.39)};
  \node[font=\footnotesize,black!60] at (8.6,1.36) {固定步长：停在 \(L\eta\sigma^2\) 上};
  \node[font=\footnotesize,black!70] at (4.55,2.72) {学习率下调};
  \node[font=\footnotesize,black!70] at (7.5,2.06) {再下调};
  \node[font=\footnotesize,black!60] at (8.7,0.16) {衰减步长：接着往下};
\end{tikzpicture}

\emph{两条曲线用同一份数据、同一个初始化，前三分之一完全重合（图中被深色覆盖）；分岔只发生在学习率被调小的那一刻。}
\end{center}

图上那两次下调各自打开了一段新的下降空间，而每段末尾又出现新的平台。这是同一个式子被用了三次：每次把 \(\eta\) 改小，\(1/(\eta T)\) 那一项变大（前期变慢），\(L\eta\sigma^2\) 那一项变小（终点更低）。工程上常见的做法是先用大步长快速走完第一项，再用小步长把第二项压下去，代价是总步数增加。

\subsection{学习率必须衰减，衰减还必须够慢}

把上面的平衡做到底，就得到衰减步长的两条要求。步长必须趋于零，否则平台一直在；但它趋于零的速度不能太快，否则参数还没走到该去的地方就动不了了。写成级数条件，这就是随机逼近里的 Robbins--Monro 条件：

\[
\sum_{t=0}^{\infty}\eta_t=\infty,\qquad \sum_{t=0}^{\infty}\eta_t^2<\infty .
\]

第一式保证总移动能力无限，参数能从任意初始位置走到任意远处；第二式保证累积的噪声方差有限，抖动能被磨掉。\(\eta_t=c/t\) 同时满足两条，\(\eta_t=c/\sqrt t\) 只满足第一条，\(\eta_t=c/t^2\) 只满足第二条——这三种情形的分界，正是 \(p\)-级数在 \(p=1\) 处的那道界，见\hyperref[sec:02-Calculus/05-Sequences-and-Series/03]{``正项级数的比较''}。同一组条件在强化学习的表格型算法里原样出现，见\hyperref[sec:16-Control-and-Reinforcement-Learning/02-Reinforcement-Learning/03]{``Monte Carlo 与随机逼近把回报变成增量更新''}。

要说清楚的是，这两条只是必要的部件，不构成收敛定理。完整结论还要求条件均值正确、高阶矩受控、迭代序列有界，以及对应的平均动力系统本身稳定。现代深度学习的日程表大多不严格满足 Robbins--Monro——余弦退火在有限步内衰减到零，warmup 阶段步长还在上升，梯度裁剪把大梯度截断后严格来说已经有偏。这些偏离不意味着算法失效，只意味着援引定理时要说明偏离在哪一条上。

另一种做法是不动步长，改动读出的参数：对轨迹上的参数做尾部平均或指数移动平均。噪声在平台上是围绕某个中心来回抖的，平均把它抵消掉一部分，等效于降了步长而不损失前期速度。这也提供了一个便宜的诊断——如果参数移动平均对应的损失明显低于瞬时损失，说明你正在平台上绕圈，而不是真的收敛了。

\subsection{batch size 只改动方差这一头}

若 batch 内样本条件独立、单样本梯度方差有界，大小为 \(B\) 的平均梯度把 \(\sigma^2\) 降到约 \(\sigma^2/B\)，平台高度随之降到 \(L\eta\sigma^2/B\)。这个式子解释了一条广为流传的经验规则：\(B\) 与 \(\eta\) 同比例放大时平台高度不变，于是大 batch 训练常配大学习率。但它是这个近似式的推论，不是定理，实际能放大的范围有上限，超过之后训练会不稳。

同样要注意，\(1/B\) 这条规律依赖抽样独立。不放回抽样、数据增强、序列内的相关性、分布式训练里的延迟更新都会改动它。batch normalization 更彻底：它让单样本的损失依赖同批的其他样本，``单样本梯度的独立平均''这个模型从一开始就不成立，见\hyperref[sec:14-Deep-Learning/03-Optimization-and-Training/03]{``归一化与 Dropout 改变训练时的随机系统''}。

固定的样本访问预算下，\(B\) 变大意味着更新次数变少。比较两个配置时看``每步梯度多准''没有意义，要看相同数据访问量或相同墙钟时间下走了多远。

\subsection{平台上的噪声不全是坏事}

上面一直把噪声当成需要消除的东西，这在凸问题上大体成立，在非凸问题上要打折扣。上一篇提到，随机扰动能帮算法从鞍点和浅盆地里出来；固定步长 SGD 的那个平台，同时也意味着算法保有在参数空间里游走的能力，不会过早锁死在第一个碰到的驻点上。经验上，先用较大的学习率跑足够久，再衰减，往往比一开始就用小学习率得到更好的最终结果——这个观察相当稳健，但目前的理论解释还不统一，宜当作实践经验而非结论使用。

需要严格区分的是两种随机性。mini-batch 方差描述的是在给定训练集上估计梯度的噪声，它属于优化层面，加大 batch 或衰减学习率就能压小。训练集本身相对于总体的抽样波动属于另一个层次，加大 batch 对它毫无作用。损失曲线的短期抖动、不同随机种子之间的差异、以及测试指标的不确定性，来自不同的随机源，得分别诊断——这条线索是本单元最后一篇的主题。

在此之前还有一个更近的问题。上面所有分析都假设更新方向就是（带噪的）负梯度方向。实际训练里几乎没人这么用：动量、Adam、RMSProp 都在梯度和更新之间插了一层加工。那层加工改变了什么，下一篇处理。
